% !PRINTED 196
\weberpage{196}

As to the degree with respect to $\varkappa^2$ it can further be concluded that
the coefficients of $x^{2\nu}$ do not exceed the degree $\nu$ with respect to
$\varkappa^2$. For if this rule is correct for $n$ and $n+1$, its correctness
follows for $2n$ and $2n+1$, and for $n=1$, $n=2$ it holds.

From the fundamental equations between the three elliptic functions:
\[
1 = \mathrm{cn}^2 v + \mathrm{sn}^2 v = \mathrm{dn}^2 v + \varkappa^2\,\mathrm{sn}^2 x % [sic]
\]
there follow further the relations:
\weq{15}{
\begin{array}{l}
D^2 = B^2 + x^2y^2z^2A^2 = C^2 + \varkappa^2x^2y^2z^2A^2,\quad n\text{ even,}\\
D^2 = y^2B^2 + x^2A^2 = z^2C^2 + \varkappa^2x^2A^2,\quad n\text{ odd,}
\end{array}}
with whose help one can, by (14), express $B_{2n}$ and $C_{2n}$ through $A_n$
and $D_n$ alone as follows:
% [sic] The German source prints (14), but the displayed recurrences for
% B_{2n} and C_{2n} are in (13); (14) concerns the odd indices 2n+1.
\weq{16}{
\begin{array}{r@{\;}c@{\;}l@{\qquad}l}
B_{2n} &=& D_n^4 - 2x^2y^2z^2A_n^2D_n^2 + \varkappa^2x^4y^4z^4A_n^4 & \\
C_{2n} &=& D_n^4 - 2\varkappa^2x^2y^2z^2A_n^2D_n^2 + \varkappa^2x^4y^4z^4A_n^4 & \left.\right\}\ n\text{ even,}\\[2mm]
B_{2n} &=& D_n^4 - 2x^2A_n^2D_n^2 + \varkappa^2x^4A_n^4 & \\
C_{2n} &=& D_n^4 - 2\varkappa^2x^2A_n^2D_n^2 + \varkappa^2x^4A_n^4 & \left.\right\}\ n\text{ odd.}
\end{array}}

\begin{center}
\textbf{\S\ 58. \gesperrt{Multiplication of the Function} $\wp(u)$.}
\end{center}

The theory of multiplication takes a very elegant form for the
\gesperrt{Weierstrass} function $\wp(u)$ with the two periods $\omega_1, \omega_2$.

Let us consider the doubly periodic function
\weq{1}{\wp(nu) - \wp(u),}
in which $n$ may be an arbitrary positive integer. This function becomes
infinite for $u=0$, and indeed in such a way that
\weq{2}{\wp(nu) - \wp(u) + \frac{n^2-1}{n^2}\,\frac{1}{u^2}}
remains finite for $u=0$. But besides this it becomes infinite for all values
$u^*$ of $u$ which satisfy the condition
\[
n\,u^* \equiv 0 \pmod{\omega_1,\omega_2}
\]
and which are therefore of the form:
\weq{3}{u^* = \frac{\nu_1\omega_1 + \nu_2\omega_2}{n},}
in which $\nu_1, \nu_2$ denote arbitrary integers. We obtain all incongruent
values contained in (3) if we let $\nu_1$ and $\nu_2$ each run through a
complete system of residues with respect to the modulus $n$.
